\uuid{YCt1}
\exo7id{4211}
\titre{exo7 4211}
\auteur{quercia}
\organisation{exo7}
\datecreate{2010-03-11}
\isIndication{false}
\isCorrection{true}
\chapitre{Equation différentielle}
\sousChapitre{Equations aux dérivées partielles}
\cpge{mpsi-mp2i.fonctions-deux-variables.b, pcsi.fonctions-deux-variables.b, ptsi.fonctions-deux-variables.b, mpsi-mp2i.fonctions-deux-variables.c, pcsi.fonctions-deux-variables.c, ptsi.fonctions-deux-variables.c, mpsi-mp2i.primitives-equations-differentielles.c, pcsi.techniques-calcul-primitives.c, ptsi.techniques-calcul-primitives.c, bcpst1.fonctions-deux-variables.c, bcpst1.fonctions-deux-variables.d, bcpst1.equations-differentielles-lineaires.b, mp-mpi.calcul-differentiel-optimisation.a}
\module{Analyse}
\niveau{L2}
\difficulte{}

\contenu{
\question{
Trouver les applications $f : {(\R^{+*})^2} \to \R$ de classe $\mathcal{C}^2$
vérifiant : $x^2\frac{\partial^2 f}{\partial x^2} + 2xy\frac{\partial^2 f}{\partial x \partial y} + y^2\frac{\partial^2 f}{\partial y^2} = 0$.
On utilisera le changement de variables : $u = xy$, $v = \frac xy$.
}
\reponse{
$2u\frac{\partial^2 g}{\partial u^2} + \frac{\partial g}{\partial u} = 0  \Rightarrow 
         f(x,y) = A\Bigl(\frac xy\Bigr)\sqrt{xy} + B\Bigl(\frac xy\Bigr)$.
}
}
